\section{Methodology Adopted}

\subsection{Investigative Techniques}
Five techniques are used: literature review, code review, component tests,
public-data experiments, and controlled field tests. Each gives different
evidence. Papers give accepted methods and known limits. Software tests check
repeatable behaviour. Field tests measure the complete system in its actual
setting.

\subsubsection{Literature and Dataset Investigation}
The literature review covers tracking, football identity, tactical analysis,
wearable sensing, and workload modelling. SoccerNet-Tracking and SportsMOT show
common sports-tracking problems \cite{cioppa2022soccernet,cui2023sportsmot}.
ByteTrack gives the starting association method \cite{zhang2022bytetrack}.
Tactical and wearable papers help select features and quality checks
\cite{honda2022pass,teixeira2025tactical,pollreisz2022wearable}.

The current SoccerNet conversion keeps ball, goalkeeper, player, and referee
classes. Its metadata records 57 training sequences and 49 test sequences from
the available source layout. A sequence is kept within one split so adjacent
frames cannot appear in both development and test data. Campus footage will be
split by recording session for the same reason.

\subsubsection{Code and Component Investigation}
The repository was checked from firmware to dashboard. The vision path contains
detection, ByteTrack association, appearance matching, team assignment, jersey
voting, pitch mapping, trajectory, pass, shot, substitution, and tactical
modules. The data path contains sensor reading, filtering, timestamp mapping,
relay ingestion, player-device binding, local API, and training-data export.

Automated tests cover packet parsing, GPS and PPG rules, relay order and replay,
player-device tags, pipeline state, training data, and API behaviour. Frontend
lint and production build checks are also used. The tested
\texttt{wearable\_stream} firmware is compiled with Arduino CLI in the
\texttt{soccer} conda environment using the official Espressif board package.
Compilation checks source and libraries;
physical accuracy is tested separately.

\subsubsection{Accuracy Investigation}
Accuracy is split into smaller questions. Detection is reported per class with
precision, recall, AP$_{50}$, and AP$_{50:95}$. Tracking is reported with HOTA,
IDF1, and identity switches. Team and jersey identity are checked against a
reviewed roster, while ball tracking also reports gaps and false recovery.

Pitch mapping uses marked reference points and known field distances. The
report will include point error in metres and repeat tests after camera setup.
Synchronisation uses visible LED events together with an IMU tap. At least 20
events will be aligned across video and wearable records, and the median, 95th
percentile, and largest absolute error will be reported.

BPM and estimated SpO$_2$ will be compared with a reference device during rest,
walking, and a controlled movement drill. The result will include error,
coverage, and the number of windows rejected by the signal gate. IMU axes will
be checked at rest and with known orientation. GPS will be checked on a marked
route with fix state and dropout time included.

Model evaluation is done only after the dataset gate is met. Player-session
records are split by time and player group where possible. A simple rule or
linear model is used as the baseline. Classification results will include
ROC-AUC, average precision, Brier score, and a reliability plot, with class
counts reported beside them.

\subsection{Proposed Solution}
The proposed solution has four working paths: vision, wearable, fusion, and
delivery. They meet at one typed event format. Each event keeps its session,
source, entity, source time, mapped time, receive time, sequence, payload, and
quality state.

\subsubsection{Vision Processing}
Frames are decoded with a monotonic frame time. The detector returns football
objects, and ByteTrack joins player detections across frames. A homography maps
tracks to the pitch. Team colour, appearance, jersey votes, and manual review
are kept as separate identity evidence. The ball has its own short recovery
rules because it is smaller and faster than a player.

\subsubsection{Wearable Acquisition}
The ESP32-S3 reads the MPU-6050, MAX30102, and NEO-6M. The main firmware sends
timestamped raw batches to the relay with a device ID, boot ID, and source
sequence. Sampling and HTTPS use separate tasks, with a fixed queue between
them. Wi-Fi credentials, relay token, CA chain, and device tags are added by the
provisioning path in an ignored generated header. A sensor failure is reported
in status events while other valid readings continue.

The host converts units and checks plausible ranges. PPG windows are checked
for contact, motion, signal size, and peak quality. GPS keeps its fix state and
age. The IMU keeps raw axes as well as filtered magnitude so workload formulas
can be checked later.

\subsubsection{Synchronisation and Fusion}
SNTP anchors the wearable clock to UTC. Each sample keeps source time, relay
receive time, and sequence. The relay orders accepted events, and a bounded
window handles a short network loss. A wearable value is joined to a track only
when the player-device binding is valid and the time error is inside the set
limit.

The fusion result never removes its source details. An identity correction is
saved as a new record with operator and time. Earlier features can then be
calculated again without changing the raw capture.

\subsubsection{Analytics and Delivery}
Simple features are calculated first: distance, speed, acceleration load,
high-intensity efforts, team centre, width, length, compactness, line spacing,
and pressure distance. Event rules point to the positions and time that caused
them. Later predictive models use reviewed player-session features and keep the
model version with every result.

FastAPI provides commands and snapshots, while WebSocket carries live events
to the dashboard. The required VPS relay receives validated raw wearable events
through \nolinkurl{cpg44.nivaspms.com}. It assigns a relay sequence and receive
time. The local sensor processor reads the live stream and can request missed
events from the fixed 4 MiB replay window. The relay has no volume, database,
video, or session archive.

\subsection{Work Breakdown Structure}
Table~\ref{tab:wbs} divides the work into products that can be shown and tested.
Each module gives evidence needed by the next module.

\begin{longtable}{@{}p{0.45in}p{1.2in}p{2.0in}p{2.0in}@{}}
\caption{Work breakdown structure}\label{tab:wbs}\\
\toprule
\textbf{WBS} & \textbf{Module} & \textbf{Main work} & \textbf{Evidence or product}\\
\midrule
\endfirsthead
\toprule
\textbf{WBS} & \textbf{Module} & \textbf{Main work} & \textbf{Evidence or product}\\
\midrule
\endhead
1.0 & Data & Convert, label, split, review & Dataset note and label report\\
2.0 & Vision & Detect, track, identify, recover ball & Tracks, metrics, error clips\\
3.0 & Calibration & Map pitch and check known points & Calibration file and error table\\
4.0 & Wearable & Wire, sample, provision, transmit & Firmware build and packet capture\\
5.0 & Sync & Anchor time, reorder, bind player & Tap-test report and fused records\\
6.0 & Analytics & Movement, workload, team geometry & Feature definitions and event cards\\
7.0 & Learning & Gate data, train baseline, evaluate & Repeatable run and model report\\
8.0 & Application & API, dashboard, correction flow & Working live application\\
9.0 & Relay & Domain route, token, cache, replay & Health and reconnect test\\
10.0 & Evaluation & Component and field protocols & Checked results and final evidence\\
\bottomrule
\end{longtable}

\subsection{Tools and Technology}
\begin{table}[H]
\centering
\caption{Tools used in the project}
\label{tab:tools}
\small
\begin{tabularx}{\textwidth}{p{1.65in}X}
\toprule
\textbf{Tool or platform} & \textbf{Use in the project}\\
\midrule
Python and conda & Vision, sensor processing, analytics, API, tests, and fixed environment\\
PyTorch and Ultralytics & Detector training and inference\\
OpenCV & Video input, pitch geometry, calibration, and review images\\
ByteTrack & Online multi-object association\\
FastAPI and WebSocket & Session API and live event transport\\
React, TypeScript, Vite & Responsive dashboard and production build\\
ESP32-S3 Arduino core & Wearable firmware, HTTPS, I2C, GPS serial, and timing\\
Arduino CLI & Repeatable firmware setup and compilation\\
Cloudflare, Caddy, Docker & Required TLS domain route and memory-only VPS relay\\
PlantUML & Editable UML and architecture diagrams\\
LaTeX and Tectonic & Report and 15-slide presentation build\\
Git & Source history and review\\
\bottomrule
\end{tabularx}
\end{table}

These tools fit the available PC and the project skill level. Python and
TypeScript keep the data format easy to inspect. PlantUML keeps diagrams easy to
change after mentor feedback, and the same rendered files are used in the
report and slides.
